\chapter{The executable kernel}
\label{chap:kernel}

This chapter describes \emph{the code that is being verified}: the Lean 4 re-implementation of
Lean's C++ kernel (\texttt{src/kernel/*.cpp}) that the \texttt{Lean4Lean/Verify/} layer reasons
about. Eighteen files and 4040 lines are in scope: the type checker
\srcloc{Lean4Lean/TypeChecker.lean}{1} (979 lines), the inductive-block adder
\srcloc{Lean4Lean/Inductive/Add.lean}{1} (776) and its reduction companion
\srcloc{Lean4Lean/Inductive/Reduce.lean}{1} (117), the primitive recognizer
\srcloc{Lean4Lean/Primitive.lean}{1} (611), universe levels
\srcloc{Lean4Lean/Level.lean}{1} (363), the replay driver
\srcloc{Lean4Lean/Replay.lean}{1} (325) with \srcloc{Main.lean}{1} (148), declaration admission
\srcloc{Lean4Lean/Environment.lean}{1} (129) over the environment utilities
\srcloc{Lean4Lean/Environment/Basic.lean}{1} (138), quotients
\srcloc{Lean4Lean/Quot.lean}{1} (121), and the small support modules
\srcloc{Lean4Lean/Expr.lean}{1} (105), \srcloc{Lean4Lean/EquivManager.lean}{1} (64),
\srcloc{Lean4Lean/Instantiate.lean}{1} (33), \srcloc{Lean4Lean/FuelConfig.lean}{1} (31),
\srcloc{Lean4Lean/ForEachExprV.lean}{1} (30), \srcloc{Lean4Lean/Declaration.lean}{1} (27),
\srcloc{Lean4Lean/PtrEq.lean}{1} (22) and \srcloc{Lean4Lean/LocalContext.lean}{1} (21).

Notation follows the thesis. $\Gamma \vdash e : \alpha$ is the typing judgment,
$\Gamma \vdash e \equiv e'$ the ideal definitional equality, $\Gamma \vdash e \Leftrightarrow e'$
the algorithmic equality that the kernel actually decides, and $e \rightsquigarrow e'$ head
reduction. The correspondence with Carneiro's \emph{The Type Theory of Lean} is close but
incomplete in both directions. The file \texttt{axioms.tex} (\S2) supplies the rules that
\texttt{inferType'}, \texttt{isDefEqCore'}, \texttt{checkPositivity},
\texttt{isLargeEliminator}, \texttt{mkRecInfos}, \texttt{mkRecRules},
\texttt{inductiveReduceRecCore}, \texttt{toCtorWhenK} and \texttt{addQuot} implement; but three
substantial pieces of the executable kernel have no thesis counterpart at all, namely
\texttt{Expr.proj} with \texttt{inferProj}/\texttt{reduceProjCore}/\texttt{toCtorWhenStruct},
nested inductive elimination (\texttt{ElimNestedInductive}), and the primitive \texttt{Nat}
arithmetic that \texttt{Primitive.checkDef} exists to justify. The first of these is exactly the
gap the \texttt{trproj} branch fills on the model side (Chapter \ref{chap:proj}).

Two design decisions dominate the code. First, nothing in the type checker is \texttt{partial}:
the four mutually recursive operations are bundled in a \texttt{Methods} record and the knot is
tied at a finite depth by \texttt{Methods.withFuel}, so every definition has usable equations for
the verification layer. The price is a real behavioural divergence, since the kernel can report
\texttt{.deepRecursion} or \texttt{.deterministicTimeout} where the C++ kernel would keep going.
Second, several definitions are shaped by their proofs rather than by the algorithm:
\texttt{reduceProj} is written with an explicit bind so that \texttt{reduceProj.WF} can see
through it, and the contributed split of \texttt{inductiveReduceRecCore} exists for the same
reason. The repository documents fourteen deliberate divergences from the C++ kernel in
\texttt{divergences.md}, which should be read as part of the deliverable.

Of the 4040 lines, exactly 22 are contributed: 21 on \texttt{trproj} in
\srcloc{Lean4Lean/Inductive/Reduce.lean}{75} (the extraction of the pure $\iota$ step into
\texttt{inductiveReduceRecCore}, plus the one call site at line 115) and 1 on \texttt{iota} in
\srcloc{Lean4Lean/Quot.lean}{24} (rejecting an \texttt{unsafe} \texttt{Eq} during quotient
initialization). Everything else in this chapter is Mario Carneiro's upstream code. The ratio is
the headline finding for this chapter, and it is a favourable one: both changes are small, both
are motivated by a concrete proof obligation rather than by taste, both are recorded in the
commit message \emph{and} in \texttt{divergences.md}, and the \texttt{iota} line discharges a
global assumption (\texttt{Environment.EqSafe}) instead of adding one. The \texttt{trproj} split
is behaviour-preserving: the extracted body is the master body verbatim, with the caller's
\texttt{let majorIdx} recomputed as \texttt{rval.getMajorIdx} at the two use sites. The only
criticisms are cosmetic (a duplicated docstring and a misleading \texttt{section} placement) and
are recorded in the review notes.

\section{Expressions, contexts and small helpers}

\begin{definition}[Delta-reducible value of a constant, \texttt{deltaValue?}]
\label{def:kernel-delta-value}
\lean{Lean.ConstantInfo.deltaValue?}
\leanok
The value a constant unfolds to under $\delta$-reduction: \texttt{defnInfo} and \texttt{thmInfo}
have one, every other \texttt{ConstantInfo} does not. This is deliberately not Lean's
\texttt{ConstantInfo.value?}, whose meaning changed in lean4\#12973 so that it excludes theorems,
while the C++ \texttt{constant\_info::has\_value()} was left untouched and the kernel therefore
still unfolds theorems.
\thesisref{axioms.tex \S2.5, the $\delta$ rule $c_{\bar\ell} \rightsquigarrow v_{\bar\ell}(c)$}
\end{definition}

\begin{definition}[Reducibility-hint ordering, \texttt{lt'}]
\label{def:kernel-reducibility-lt}
\lean{Lean.ReducibilityHints.lt'}
\leanok
The strict order on \texttt{ReducibilityHints} that lazy delta reduction uses to choose which of
two $\delta$-reducible heads to unfold: nothing is below \texttt{opaque}, \texttt{abbrev} is below
nothing, and \texttt{regular $d_1$} is below \texttt{regular $d_2$} iff $d_1 < d_2$. No thesis
counterpart; it is a heuristic of the algorithm, not a rule of the theory.
\reviewnote{The single marker \texttt{-- lean4\#2750} is the only record that the C++ hint
comparison is known to differ subtly from this one.}
\end{definition}

\begin{definition}[Expression constructors and recognizers]
\label{family:kernel-expr-helpers}
\lean{Lean.Expr.prop, Lean.Expr.arrow, Lean.Expr.lam0, Lean.Expr.natZero, Lean.Expr.natSucc, Lean.Expr.isConstructorApp?'}
\leanok
Short names for expressions the kernel builds constantly: $\mathsf{Prop} = \mathsf{Sort}\;0$, the
non-dependent arrow, an anonymous lambda, and the constants \texttt{Nat.zero} and
\texttt{Nat.succ}. \texttt{isConstructorApp?'} returns the head constant of an application when
that constant is a constructor of the given environment, and \texttt{none} otherwise.
\end{definition}

\begin{definition}[Monadic expression replacement, \texttt{replaceM}]
\label{def:kernel-expr-replace}
\lean{Lean.Expr.replaceM, Lean.Expr.replaceNoCacheT, Lean.Expr.ReplaceImpl.replaceUnsafe', Lean.Expr.ReplaceImpl.replaceUnsafeT}
\leanok
\texttt{replaceM f? e} rewrites $e$ top-down by the partial function \texttt{f?}, recursing into
subterms wherever \texttt{f?} declines. The reference implementation \texttt{replaceNoCacheT} is
structurally recursive; a pointer-keyed memoizing \texttt{unsafe} version is installed over it
with \texttt{@[implemented\_by]}.
\reviewnote{\texttt{replaceM} is itself \texttt{partial} even though \texttt{replaceNoCacheT} is
not, so it has no equation lemmas, and the cached implementation is trusted. A \texttt{TODO} at
\srcloc{Lean4Lean/Expr.lean}{44} acknowledges this.}
\end{definition}

\begin{definition}[Literal expansion to constructor applications]
\label{def:kernel-lit-to-ctor}
\lean{Lean.Expr.natLitToConstructor, Lean.Expr.strLitToConstructor, Lean.Literal.toConstructor, Lean.Literal.typeName}
\leanok
Expansion of the kernel's compact literals into constructor applications: $0 \mapsto
\mathsf{Nat.zero}$ and $n+1 \mapsto \mathsf{Nat.succ}\;n$ with the predecessor kept as a literal,
and a string literal into \texttt{String.ofList} applied to an explicit \texttt{List Char} of
\texttt{Char.ofNat} applications. \texttt{typeName} names the type of a literal. This is the
bridge that lets the $\iota$ rule fire on literals, and hence the reason the primitive status of
\texttt{Nat} is a soundness obligation (Definition \ref{def:kernel-primitive-checkinductive}).
\end{definition}

\begin{definition}[Substitution-free beta reduction, \texttt{cheapBetaReduce}]
\label{def:kernel-cheap-beta}
\lean{Lean.Expr.cheapBetaReduce}
\leanok
Beta-reduces $(\lambda x_1\dots x_n.\;b)\,a_1\dots a_m$ in exactly the two cases that need no
substitution: to $b\,a_{n+1}\dots a_m$ when $b$ has no loose bound variables, and to
$a_i\,a_{n+1}\dots a_m$ when $b$ is the bound variable $x_i$. Any other term is returned
unchanged. It is used to tidy the inferred types produced by \texttt{inferLambda} and
\texttt{inferLet}.
\thesisref{axioms.tex \S2.2, the $\beta$ rule, restricted to substitution-free instances}
\reviewnote{The $x_i$ branch is guarded by \texttt{assert! n < i} at
\srcloc{Lean4Lean/Instantiate.lean}{24}; the invariant is not established here.}
\end{definition}

\begin{definition}[Scoped local declarations]
\label{def:kernel-monad-localnamegen}
\lean{Lean4Lean.MonadLocalNameGenerator, Lean4Lean.withLocalDecl, Lean4Lean.withLetDecl}
\leanok
\texttt{MonadLocalNameGenerator} abstracts the supply of fresh \texttt{FVarId}s;
\texttt{withLocalDecl} and \texttt{withLetDecl} extend the ambient \texttt{LocalContext} with a
fresh free variable of a given type (respectively, with a value) for the duration of a
continuation. Every telescope walk in the kernel uses these, so they are how the context
extension $\Gamma, x{:}\alpha$ of the thesis becomes an operation on a concrete context.
\thesisref{axioms.tex \S2.1, context extension $\Gamma, x{:}\alpha$}
\end{definition}

\begin{definition}[Cached expression traversal without opaques, \texttt{forEachV}]
\label{def:kernel-foreach-exprv}
\lean{Lean.Expr.forEachV, Lean.Expr.forEachV', Lean.ForEachExprV.visit}
\leanok
A copy of \texttt{Expr.forEach} built on \texttt{MonadStateCacheT} rather than
\texttt{StateRefT}, so the traversal is a plain definition instead of an \texttt{opaque} and can
therefore be unfolded in proofs.
\reviewnote{The module is imported by \texttt{TypeChecker.lean} but neither \texttt{forEachV} nor
\texttt{forEachV'} has a call site anywhere under \texttt{Lean4Lean/}; this looks like dead code
left from an earlier version of the expression scope check.}
\end{definition}

\section{Universe levels}

\begin{definition}[Undefined level parameter detection, \texttt{getUndefParam}]
\label{def:kernel-level-getundefparam}
\lean{Lean.Level.getUndefParam, Lean.Level.forEach}
\leanok
\texttt{getUndefParam l ps} returns some parameter occurring in $\ell$ but not in the
declaration's list $\bar u$ of universe variables, by a short-circuiting \texttt{forEach} over the
level. \texttt{checkLevel} (Definition \ref{def:kernel-infer-constant}) uses it to enforce the
thesis's side condition that the universe variables of a declaration are contained in $\bar u$.
\thesisref{axioms.tex \S2.5, ``the universe variables in $\alpha$ and $\ell$ are contained in $\bar u$''}
\end{definition}

\begin{definition}[Canonical form for universe levels, \texttt{Normalize.normalize}]
\label{def:kernel-level-normalize}
\lean{Lean.Level.Normalize.normalize, Lean.Level.Normalize.normalizeAux, Lean.Level.Normalize.NormLevel, Lean.Level.Normalize.Node, Lean.Level.Normalize.VarNode, Lean.Level.Normalize.NormLevel.subsumption, Lean.Level.Normalize.NormLevel.toTree, Lean.Level.Normalize.Tree.reify, Lean.Level.normalize'}
\leanok
A canonical form for the level algebra $\ell ::= u \mid 0 \mid \ell+1 \mid \max \mid \mathrm{imax}$,
following G\'eran's ``A Canonical Form for Universe Levels in Impredicative Type Theory''. A level
becomes a map from $\mathrm{imax}$ condition sets (sorted lists of parameter names) to a
\texttt{Node} recording the largest constant and the largest offset of each parameter;
\texttt{subsumption} deletes dominated sublevels, and \texttt{toTree}/\texttt{Tree.reify} read a
canonical \texttt{Level} back out, choosing the lexicographically least admissible
$\mathrm{imax}$ chain so that the output depends only on the canonical sublevels.
\thesisref{axioms.tex \S2.2, the level syntax and the judgment $\ell \le \ell' + n$}
\reviewnote{Substantial new algorithmic content, but its correctness is not established in this
chapter's scope: Chapter \ref{chap:levels} proves it, in full and sorry-free, as
\texttt{Lean.Level.Normalize.normalize\_complete} and \texttt{Lean.Level.isEquiv'\_complete}
($\texttt{isEquiv'}\;u\;v \leftrightarrow u \approx v$ against the semantic model
\texttt{VLevel}), so this is not the unverified new component it might appear to be in isolation.}
\end{definition}

\begin{definition}[Sublevel comparison on canonical forms, \texttt{NormLevel.le}]
\label{def:kernel-level-le}
\lean{Lean.Level.Normalize.NormLevel.le, Lean.Level.Normalize.NormLevel.addable, Lean.Level.Normalize.NormLevel.feasible, Lean.Level.Normalize.NormLevel.lexChain}
\leanok
\uses{def:kernel-level-normalize}
Decides $\ell_1 \le \ell_2$ on canonical forms: every sublevel of $\ell_1$ must be dominated by
sublevels of $\ell_2$, where each entry of $\ell_2$ discharges what it can of an outstanding node
of $\ell_1$ and the node counts as dominated once nothing is left.
\texttt{addable}/\texttt{feasible}/\texttt{lexChain} build the lexicographically least admissible
condition chain used by reification.
\thesisref{axioms.tex \S2.2, the algorithmic inequality $\ell \le \ell' + n$}
\reviewnote{Cited in this docstring only as ``Theorem 39 of the paper'', with no formal
bibliography entry, though the paper itself (Yoan G\'eran, ``A Canonical Form for Universe Levels
in Impredicative Type Theory'') is named with a working URL earlier in the same file, at
\srcloc{Lean4Lean/Level.lean}{27}. In any case the claim is not merely cited on trust: Chapter
\ref{chap:levels} reproves it from scratch in Lean, sorry-free, as
\texttt{Lean.Level.Normalize.NormLevel.le\_complete} and \texttt{Lean.Level.geq'\_complete}.}
\end{definition}

\begin{definition}[Complete level equality and comparison, \texttt{isEquiv'}]
\label{def:kernel-level-isequiv}
\lean{Lean.Level.isEquiv', Lean.Level.isEquivList, Lean.Level.geq'}
\leanok
\uses{def:kernel-level-normalize, def:kernel-level-le}
$\ell_1 \equiv \ell_2$ is decided by running core Lean's sound-but-incomplete \texttt{isEquiv}
first and falling back to equality of canonical forms; $\ell_1 \ge \ell_2$ likewise falls back to
\texttt{NormLevel.le}. \texttt{isEquivList} lifts equality pointwise to the level arguments of a
constant. This is a documented divergence in the \emph{accept-more} direction: lean4lean decides
level equality and $\ge$ for strictly more pairs than either the standard library or the C++
kernel, so it accepts declarations the kernel rejects for level incompleteness and never the
converse.
\thesisref{axioms.tex \S2.2, $\ell \le \ell'$ and $\ell' \le \ell$ imply $\ell \equiv \ell'$}
\reviewnote{The docstring quotes a measurement (261k comparisons over Lean, Std and Batteries;
roughly 20 times cheaper than normalizing) that no script in the repository reproduces.}
\end{definition}

\section{The kernel environment}

\begin{definition}[Reserved primitive names, \texttt{primitives}]
\label{def:kernel-primitives}
\lean{Lean.Kernel.Environment.primitives}
\leanok
The name set with special kernel support: \texttt{Bool} and its two constructors, \texttt{Nat}
with \texttt{zero}/\texttt{succ} and its arithmetic and bitwise operations, \texttt{String.ofList}
and \texttt{Char.ofNat}. A declaration may claim one of these names only if
\texttt{Primitive.checkDef} or \texttt{Primitive.checkInductive} has certified that it has the
expected shape.
\end{definition}

\begin{definition}[Environment lookup and pre-addition hygiene checks]
\label{family:kernel-env-checks}
\lean{Lean.Kernel.Environment.contains, Lean.Kernel.Environment.get, Lean.Kernel.Environment.checkDuplicatedUnivParams, Lean.Kernel.Environment.checkNoMVar, Lean.Kernel.Environment.checkNoFVar, Lean.Kernel.Environment.checkNoMVarNoFVar, Lean.Kernel.Environment.checkName}
\leanok
\uses{def:kernel-primitives}
Lookup of a constant by name (\texttt{get}, failing with \texttt{unknownConstant}), rejection of a
repeated universe parameter in a declaration's level parameter list, rejection of metavariables
and free variables in a declaration's type and value, and \texttt{checkName}, which rejects a name
already present in the environment and, unless \texttt{allowPrimitive} is set, any reserved
primitive name.
\thesisref{axioms.tex \S2.5, the admissibility side conditions on $\mathsf{constant}\;c_{\bar u}:\alpha$ and $\mathsf{def}\;c_{\bar u}:\alpha := e$}
\end{definition}

\begin{definition}[Non-recursive structure predicate, \texttt{isNonRecStructure}]
\label{def:kernel-is-nonrec-structure}
\lean{Lean.Kernel.Environment.isNonRecStructure}
\leanok
True exactly when the name resolves to an \texttt{inductInfo} with \texttt{isRec := false}, a
single constructor and no indices. This delimits the class of types for which the kernel supports
structure $\eta$ and \texttt{Expr.proj}, and the docstring records that it must stay in sync with
the C++ \texttt{is\_non\_rec\_structure()}.
\reviewnote{The thesis's type theory has neither projections nor structure $\eta$, so this
predicate delimits precisely the extension the \texttt{trproj} model work has to cover
(Chapter \ref{chap:proj}).}
\end{definition}

\begin{definition}[Empty and imported kernel environments]
\label{def:kernel-env-empty-import}
\lean{Lean.Kernel.Environment.empty, Lean.Kernel.Environment.finalizeImport, Lean.Kernel.Environment.throwAlreadyImported}
\leanok
The two constructors the checker starts from. \texttt{empty} builds a fresh environment for a main
module; \texttt{finalizeImport} merges the constant maps of a list of imported modules, keeping a
declaration that subsumes the one already recorded and rejecting a genuine clash. Both reach into
Lean internals through \texttt{open private}, so they are sensitive to core refactors.
\reviewnote{\srcloc{Lean4Lean/Environment/Basic.lean}{124} sets \texttt{quotInit := !imports.isEmpty}
unconditionally, that is, it \emph{assumes} any non-empty import set has already initialized
quotients. \texttt{Replay.checkQuotInit} is the compensating check, but an imported environment is
never re-checked, so \texttt{checkEqType} never runs on one.}
\end{definition}

\section{Pointer equality, fuel and the equivalence cache}

\begin{definition}[Pointer equality and its soundness axioms]
\label{axiom:kernel-ptr-eq}
\lean{Lean4Lean.ptrEqExpr, Lean4Lean.ptrEqExpr_eq, Lean4Lean.ptrEqConstantInfo, Lean4Lean.ptrEqConstantInfo_eq}
\leanok
\texttt{ptrEqExpr} and \texttt{ptrEqConstantInfo} are \texttt{opaque} Booleans implemented by
\texttt{ptrAddrUnsafe} comparison, and \texttt{ptrEqExpr\_eq}/\texttt{ptrEqConstantInfo\_eq} are
\texttt{axiom}s stating that a \texttt{true} result implies actual equality. The converse is
correctly \emph{not} assumed. The docstring explains why \texttt{withPtrEq} cannot be used: the
kernel's behaviour, not merely its speed, depends on pointer identity, since identical objects
receiving different addresses can make it reject inputs it would otherwise accept.
\axfoot{Lean4Lean.ptrEqExpr\_eq; Lean4Lean.ptrEqConstantInfo\_eq; Classical.choice}
\reviewnote{These are two of the project's genuine trusted axioms, used in only two places
(\srcloc{Lean4Lean/Verify/EquivManager.lean}{264} and
\srcloc{Lean4Lean/Verify/TypeChecker/IsDefEq.lean}{384}), so the trusted surface is small. But at
\srcloc{Lean4Lean/TypeChecker.lean}{739} \texttt{ptrEqConstantInfo dt ds} gates a \emph{different
algorithm} (congruence on arguments before unfolding), not merely a shortcut, which is exactly
what the file's own docstring warns about.}
\end{definition}

\begin{definition}[Fuel bounds for the kernel's loops, \texttt{FuelConfig}]
\label{structure:kernel-fuel-config}
\lean{Lean4Lean.FuelConfig}
\leanok
One \texttt{Nat} per bounded loop: the \texttt{whnf'} unfold loop in its ordinary
(\texttt{whnf}, default 100000) and eager (\texttt{whnfEager}, 1000000) forms,
\texttt{lazyDelta} (1000), \texttt{etaExpand} (1000), the starting fuel \texttt{recDepth} (50000)
for \texttt{Methods.withFuel}, and \texttt{inductiveFuel} (1000) shared by the structural loops of
the inductive adder. Exhaustion throws \texttt{.deterministicTimeout} or \texttt{.deepRecursion}.
This replaces the C++ kernel's \texttt{maxRecDepth} (lean4\#13956) and is recorded as a divergence.
\reviewnote{Every field's docstring names the loop it bounds, which is exemplary; the docstring's
claim that ``defaults are set so mathlib passes'' is backed by nothing in the repository.}
\end{definition}

\begin{definition}[Union-find cache of definitional equalities, \texttt{EquivManager}]
\label{def:kernel-equiv-manager}
\lean{Lean4Lean.EquivManager, Lean4Lean.EquivManager.isEquiv, Lean4Lean.EquivManager.addEquiv, Lean4Lean.EquivManager.find, Lean4Lean.EquivManager.merge, Lean4Lean.EquivManager.toNode}
\leanok
\uses{axiom:kernel-ptr-eq}
A \texttt{Batteries.UnionFind} over expressions recording pairs already found definitionally
equal, plus \texttt{isEquiv}, a structural comparison that short-circuits on pointer equality, on
a hash mismatch when \texttt{useHash} is set, and on a shared union-find root, then recurses
congruentially and merges the two roots on success. \texttt{addEquiv} records a successful
\texttt{isDefEq} verdict.
\reviewnote{Since \texttt{addEquiv} records pairs equated by the \emph{algorithmic} relation, and
that relation is not transitive, the union-find's transitive closure can claim equalities the
algorithm would not decide. This matches the C++ kernel, but it means \texttt{isEquiv} is not a
conservative optimization and the verification layer has to account for it.}
\end{definition}

\section{Quotient types}

\begin{definition}[Pure expression-building monad, \texttt{ExprBuildT}]
\label{def:kernel-exprbuildt}
\lean{Lean4Lean.ExprBuildT, Lean4Lean.ExprBuildT.run}
\leanok
\uses{def:kernel-monad-localnamegen}
\texttt{ExprBuildT m} is \texttt{ReaderT LocalContext (ReaderT NameGenerator m)}: a local context
plus a deterministic name generator, which is exactly enough to run \texttt{withLocalDecl} and
\texttt{withLetDecl} while building closed expressions without a type checker. It is used to
write the quotient constants' types as telescopes. Because \texttt{run} starts the generator from
a fixed state, the resulting expressions are literally computable, which is what lets
\srcloc{Lean4Lean/Verify/Environment/Quot.lean}{1} evaluate those telescopes.
\end{definition}

\begin{definition}[Validation of the prelude's \texttt{Eq}, \texttt{checkEqType}]
\label{def:kernel-check-eq-type}
\contrib{mixed} \srcloc{Lean4Lean/Quot.lean}{20}
\lean{Lean4Lean.checkEqType}
\leanok
\uses{family:kernel-env-checks, def:kernel-exprbuildt, def:kernel-monad-localnamegen}
Checks, before quotient initialization, that the environment's \texttt{Eq} is an inductive with
exactly one universe parameter $u$ and one constructor, that its type is
$\forall\,\{\alpha : \mathsf{Sort}\;u\},\ \alpha \to \alpha \to \mathsf{Prop}$ and the
constructor's is $\forall\,\{\alpha : \mathsf{Sort}\;u\}\,(a:\alpha),\ \mathsf{Eq}\;\alpha\;a\;a$
(compared as built terms, so up to binder names), and, added on \texttt{iota} at
\srcloc{Lean4Lean/Quot.lean}{24}, that \texttt{Eq} is not declared \texttt{unsafe}.
\thesisref{axioms.tex \S2.6.2 (sec:large\_elim), $\mathsf{eq}_a := \mu\,t:\alpha\to\mathsf{Prop}.\,(\mathsf{refl} : t\;a)$}
\axfoot{propext; Classical.choice; Quot.sound (no sorryAx)}
The \texttt{iota} line is the only change to kernel code on that branch and is a deliberate
divergence in the reject-more direction: \texttt{quot.cpp}'s \texttt{check\_eq\_type} pins the
shape of \texttt{Eq} but not its safety, the four quotient constants carry no safety flag and are
therefore safe, and \texttt{Quot.lift}'s type mentions \texttt{Eq}, so an \texttt{unsafe
inductive Eq} would leave safe constants depending on an unsafe one and the safe-level abstract
environment would have no model. It pays for itself: \texttt{addQuot.WF} and \texttt{addDecl.WF}
both lost the \texttt{Environment.EqSafe} hypothesis, which was deleted outright.
\reviewnote{The guarantee is narrower than the commit message suggests.
\texttt{Environment.addQuot} returns immediately when \texttt{env.quotInit} is already true, and
\texttt{finalizeImport} sets \texttt{quotInit} for any non-empty import set, so on an imported
environment neither the shape check nor the new safety check ever runs. The new property holds
only for environments built from scratch, where \texttt{Init.Prelude} is replayed. This is a
pre-existing gap, not one the contribution introduced, but it bounds what removing
\texttt{EqSafe} buys.}
\end{definition}

\begin{definition}[Quotient initialization, \texttt{addQuot}]
\label{def:kernel-add-quot}
\lean{Lean4Lean.Environment.addQuot}
\leanok
\uses{def:kernel-check-eq-type, family:kernel-env-checks, def:kernel-exprbuildt, def:kernel-monad-localnamegen}
Returns the environment unchanged if \texttt{quotInit} is already set; otherwise runs
\texttt{checkEqType}, checks that none of the four names is taken, and adds \texttt{Quot},
\texttt{Quot.mk}, \texttt{Quot.lift} and \texttt{Quot.ind} as \texttt{quotInfo} declarations with
types built explicitly as telescopes, marking the environment quotient-initialized.
\texttt{Quot.sound} is not added here: it is an ordinary axiom of the prelude, exactly as in the
thesis, where it is the only genuine axiom among the quotient constants.
\thesisref{axioms.tex \S2.7, $\alpha/R : \mathcal{U}_u$, $\mathrm{mk}_R$, $\mathrm{lift}_R$, $\mathrm{sound}_R$}
\end{definition}

\begin{definition}[Quotient computation rule, \texttt{quotReduceRec}]
\label{def:kernel-quot-reduce-rec}
\lean{Lean4Lean.quotReduceRec}
\leanok
Reduces $\mathrm{lift}\;\alpha\;r\;\beta\;f\;h\;(\mathrm{mk}\;r\;a)\;\bar c \rightsquigarrow
f\;a\;\bar c$ and $\mathrm{ind}\;\alpha\;r\;\beta\;p\;(\mathrm{mk}\;r\;a)\;\bar c \rightsquigarrow
p\;a\;\bar c$: it normalizes the argument at the fixed major position (index 5 for
\texttt{Quot.lift}, 4 for \texttt{Quot.ind}), requires it to be a \texttt{Quot.mk} application of
arity 3, applies the function argument (index 3 in both cases) to its last argument, and
re-applies any trailing arguments.
\thesisref{axioms.tex \S2.7, $\mathrm{lift}_R\;\beta\;f\;h\;(\mathrm{mk}_R\;a) \rightsquigarrow f\;a$}
\reviewnote{The positions 5/4 and 3/3 are magic numbers justified only by the type signatures
quoted in the surrounding docstring. Unlike the inductive $\iota$ rule, nothing here checks that
the eliminator is saturated: a partially applied \texttt{Quot.lift} simply fails to match.}
\end{definition}

\section{Recursor reduction}

\begin{definition}[Nullary constructor application for a type]
\label{def:kernel-get-first-ctor}
\lean{Lean4Lean.getFirstCtor, Lean4Lean.mkNullaryCtor}
\leanok
\texttt{getFirstCtor} returns the head of an inductive's constructor list;
\texttt{mkNullaryCtor type nparams} decomposes $I\;\bar a$ and returns $I$'s first constructor at
the same levels applied to the first \texttt{nparams} of those arguments, that is, applied to the
parameters only.
\end{definition}

\begin{definition}[K-like reduction, \texttt{toCtorWhenK}]
\label{def:kernel-to-ctor-when-k}
\lean{Lean4Lean.toCtorWhenK}
\leanok
\uses{def:kernel-get-first-ctor}
For a recursor with \texttt{rval.k} set, replaces a major premise $e$ whose type weak-head
normalizes to an application of the recursor's own inductive by the nullary constructor
application of that same type, provided the two types are definitionally equal (which pins the
indices) and no index carries a metavariable. Otherwise $e$ is returned unchanged. For example
$e : a = a$ becomes $\mathsf{Eq.refl}\;a$, which is definitionally equal to $e$ by proof
irrelevance.
\thesisref{axioms.tex \S2.6.4 (sec:iota), the K-like reduction rule}
\reviewnote{The thesis describes this step as ``proof irrelevance to change the major premise into
a constructor, followed by the iota rule'', which is exactly the factoring the \texttt{trproj}
branch made explicit in code. \texttt{reduceRecursor.WF} remains open precisely because this step
has no \texttt{IsDefEq} counterpart in the model yet.}
\end{definition}

\begin{definition}[Structure-eta expansion of a major premise, \texttt{toCtorWhenStruct}]
\label{def:kernel-to-ctor-when-struct}
\lean{Lean4Lean.toCtorWhenStruct, Lean4Lean.expandEtaStruct}
\leanok
\uses{def:kernel-get-first-ctor, def:kernel-is-nonrec-structure}
When the major's inductive is a non-recursive structure, the major is not already a constructor
application, its type normalizes to that structure, and the structure's universe
\texttt{isNeverZero}, replaces $e$ by
$\mathsf{S.mk}\;\bar a\;(e.0)\dots(e.(n{-}1))$, which is definitionally equal to $e$ by structure
$\eta$, so that the $\iota$ rule can then fire. No thesis counterpart: the thesis's system has
neither projections nor structure $\eta$.
\reviewnote{Second site of the documented \texttt{isNeverZero} versus \texttt{!isAlwaysZero}
divergence from the C++ kernel; it narrows the class of projections any \texttt{TrProj} theorem
can be about.}
\end{definition}

\begin{definition}[Recursor rule lookup, \texttt{getRecRuleFor}]
\label{def:kernel-get-rec-rule-for}
\lean{Lean4Lean.getRecRuleFor}
\leanok
Given a recursor's \texttt{RecursorVal} and a major premise, returns the \texttt{RecursorRule}
whose \texttt{ctor} field matches the head constant of the major, and \texttt{none} if the head is
not a constant or no rule matches. This is the selection of the clause $e_c$ for the constructor
$c$.
\thesisref{axioms.tex \S2.6.4 (sec:iota), selecting the minor premise $e_c$ for constructor $c$}
\end{definition}

\begin{definition}[The pure $\iota$ step, \texttt{inductiveReduceRecCore}]
\label{def:kernel-iota-core}
\contrib{trproj} \srcloc{Lean4Lean/Inductive/Reduce.lean}{75}
\lean{Lean4Lean.inductiveReduceRecCore}
\leanok
\uses{def:kernel-get-rec-rule-for}
A pure, monad-free \texttt{Option Expr}-valued function performing the $\iota$ rule on a major
premise that is already a constructor application. Given \texttt{rval}, the recursor's level
arguments \texttt{ls}, the redex's argument array \texttt{recArgs} and the major, it looks up the
rule for the major's head constructor, instantiates \texttt{rule.rhs} at \texttt{rval.levelParams
:= ls}, and applies the result to \texttt{recArgs} up to \texttt{rval.getFirstIndexIdx} (the
parameters, motives and minor premises; the indices are determined by the constructor), then to
the \emph{last} \texttt{rule.nfields} arguments of the major (its fields), then to any arguments
of the redex after position \texttt{rval.getMajorIdx}. It returns \texttt{none} if no rule
matches, if the major has fewer than \texttt{rule.nfields} arguments, or if the number of level
arguments disagrees with \texttt{rval.levelParams}.
\thesisref{axioms.tex \S2.6.4 (sec:iota), $\mathrm{rec}_P\;C\;e\;p[b]\;(c\;b) \rightsquigarrow e_c\;b\;v$}
\axfoot{propext; Quot.sound (no sorryAx)}
This is the whole of the \texttt{trproj} change to kernel code. The body is master's verbatim: the
only textual difference is that the caller's \texttt{let majorIdx := info.getMajorIdx} binding
becomes \texttt{rval.getMajorIdx} at the two use sites, which is the same value, so the split is
behaviour-preserving as the commit claims. The purpose is proof-side: \texttt{inductiveReduceRec}
is monadic only because of its K-like and structure-$\eta$ preludes, whereas the $\iota$ step
itself is pure and total, so extracting it lets
\srcloc{Lean4Lean/Verify/TypeChecker/WHNF.lean}{17} state
\texttt{inductiveReduceRecCore.WF} about a function with no monad, no environment lookup and no
well-typedness side conditions beyond its hypotheses. The commit is explicit that
\texttt{reduceRecursor.WF} therefore stays open, which is an honest scoping statement rather than
an overclaim. The extracted function is also used as a differential oracle by
\srcloc{Lean4Lean/Tests/IotaShape.lean}{151}.
\reviewnote{Three cosmetic blemishes. (a) The definition sits inside the \texttt{section} opened at
\srcloc{Lean4Lean/Inductive/Reduce.lean}{9} whose variables (\texttt{[Monad m]}, \texttt{env},
\texttt{whnf}, \texttt{inferType}, \texttt{isDefEq}) it uses none of; Lean binds only used section
variables so the arity is the intended four, but a later edit that touches \texttt{env} would
silently change the signature and break the \texttt{WF} statement. (b) The new docstring
duplicates the slicing description in the unchanged caller docstring. (c) It collapses ``no rule
for this constructor'', ``major has too few arguments'' and ``wrong number of universe
arguments'' into one \texttt{none}, which is why the \texttt{WF} theorem must take saturation as
an external hypothesis rather than deriving it; the commit message says so plainly.}
\end{definition}

\begin{definition}[Recursor reduction, \texttt{inductiveReduceRec}]
\label{def:kernel-inductive-reduce-rec}
\contrib{mixed} \srcloc{Lean4Lean/Inductive/Reduce.lean}{100}
\lean{Lean4Lean.inductiveReduceRec}
\leanok
\uses{def:kernel-to-ctor-when-k, def:kernel-to-ctor-when-struct, def:kernel-iota-core, def:kernel-lit-to-ctor}
Performs recursor reduction on a redex whose head is a recursor constant: it reads the major
premise at \texttt{info.getMajorIdx}, applies K-like reduction when \texttt{info.k}, weak-head
normalizes the major, converts a \texttt{Nat} literal by \texttt{natLitToConstructor}, a
\texttt{String} literal by normalizing \texttt{strLitToConstructor}, and anything else by
\texttt{toCtorWhenStruct}, and then delegates the $\iota$ step to
\texttt{inductiveReduceRecCore}. It returns \texttt{none} when no rule applies.
\thesisref{axioms.tex \S2.6.4 (sec:iota), the $\iota$ rule and the K-like rule}
\axfoot{propext; Classical.choice; Quot.sound (no sorryAx)}
This is a master definition whose last line (\srcloc{Lean4Lean/Inductive/Reduce.lean}{115}) was
changed on \texttt{trproj} to call the extracted core.
\reviewnote{The master docstring at \srcloc{Lean4Lean/Inductive/Reduce.lean}{91} was left
untouched and still describes the $\iota$ slicing in full, work that is now done in
\texttt{inductiveReduceRecCore} and documented there too. Trimming it to the preprocessing the
caller still performs would have completed the refactor.}
\end{definition}

\section{The type checker monad}

\begin{definition}[The type checker monad, \texttt{TypeChecker.M}]
\label{structure:kernel-tc-monad}
\lean{Lean4Lean.TypeChecker.State, Lean4Lean.TypeChecker.Context, Lean4Lean.TypeChecker.M, Lean4Lean.TypeChecker.M.run}
\leanok
\uses{structure:kernel-fuel-config, def:kernel-equiv-manager}
$\texttt{M} = \texttt{ReaderT Context (StateT State (Except Exception))}$. The context carries the
environment, the local context $\Gamma$, the \texttt{DefinitionSafety} level, the
\texttt{eagerReduce} flag, the declaration's universe parameters $\bar u$ and the fuel bounds; the
state carries the two inference caches (one per \texttt{inferOnly} mode), the \texttt{whnfCore}
and \texttt{whnf} caches, the name generator, the equivalence manager, the failure cache and the
delta-unfolding cache.
\reviewnote{Keeping \texttt{lparams} fixed for the whole run is one of the documented divergences:
the C++ kernel sets and unsets level parameters around \texttt{check}, so its \texttt{ensure\_sort}
runs without them.}
\end{definition}

\begin{definition}[Fuel-bounded mutual recursion, \texttt{Methods.withFuel}]
\label{structure:kernel-methods-fuel}
\lean{Lean4Lean.TypeChecker.Methods, Lean4Lean.TypeChecker.RecM, Lean4Lean.TypeChecker.Methods.withFuel, Lean4Lean.TypeChecker.RecM.run}
\leanok
\uses{structure:kernel-fuel-config, structure:kernel-tc-monad}
The four mutually recursive operations \texttt{isDefEqCore}, \texttt{whnfCore}, \texttt{whnf} and
\texttt{inferType} are fields of a record \texttt{Methods} threaded through $\texttt{RecM} =
\texttt{ReaderT Methods M}$. \texttt{Methods.withFuel 0} makes all four throw
\texttt{.deepRecursion}; \texttt{withFuel (n+1)} fills them with the primed implementations
applied to \texttt{withFuel n}. \texttt{RecM.run} starts the knot at \texttt{FuelConfig.recDepth}.
This is what replaces \texttt{partial def} and gives every definition in the file equations to
reason about.
\reviewnote{Central design decision of the project, and the source of a real behavioural
divergence: a declaration the C++ kernel accepts can be rejected here with
\texttt{.deepRecursion}.}
\end{definition}

\section{Weak-head normalization}

\begin{definition}[The $\iota$ step of normalization, \texttt{reduceRecursor}]
\label{def:kernel-reduce-recursor}
\lean{Lean4Lean.TypeChecker.Inner.reduceRecursor}
\leanok
\uses{def:kernel-quot-reduce-rec, def:kernel-inductive-reduce-rec}
Tries quotient reduction first, but only when the environment reports \texttt{quotInit}, and then
inductive recursor reduction, returning the first reduct found and \texttt{none} otherwise.
\thesisref{axioms.tex \S2.6.4 (sec:iota) and \S2.7}
\reviewnote{\texttt{reduceRecursor.WF} is explicitly still open in the verification layer, since
the K-like and structure-$\eta$ paths have no \texttt{IsDefEq} counterpart; the \texttt{trproj}
contribution verifies only the pure $\iota$ core beneath it.}
\end{definition}

\begin{definition}[Projection reduction, \texttt{reduceProj}]
\label{def:kernel-reduce-proj}
\lean{Lean4Lean.TypeChecker.Inner.reduceProjCore, Lean4Lean.TypeChecker.Inner.reduceProj}
\leanok
\texttt{reduceProjCore idx s} expands a string literal to a \texttt{String.ofList} application,
then, if the head is a constructor, returns the argument at position
$\texttt{numParams} + \texttt{idx}$ (and \texttt{none} if that position does not exist).
\texttt{reduceProj} first normalizes the structure argument, with \texttt{whnfCore} instead of
\texttt{whnf} when \texttt{cheapProj} is set. No thesis counterpart: this is the reduction rule
the \texttt{trproj} branch models. \texttt{reduceProj}, \texttt{whnfCore} (\ref{def:kernel-whnf-core})
and \texttt{whnf} (\ref{def:kernel-whnf}) are mutually recursive; to keep the dependency graph
acyclic only the back-edges from those later nodes to this one are recorded via \texttt{\textbackslash uses}.
\reviewnote{The comment at \srcloc{Lean4Lean/TypeChecker.lean}{356} states that the explicit
\texttt{>>=} is written that way because \texttt{reduceProj.WF} ``cannot see through'' the
\texttt{do}-elaborator's bind lifting. That is the implementation being bent to suit its proof;
defensible in a verification-first project, but worth naming.}
\end{definition}

\begin{definition}[Delta reduction, \texttt{unfoldDefinition}]
\label{def:kernel-delta-unfold}
\lean{Lean4Lean.TypeChecker.Inner.isDelta, Lean4Lean.TypeChecker.Inner.instantiateDeltaValue, Lean4Lean.TypeChecker.Inner.unfoldDefinitionCore, Lean4Lean.TypeChecker.Inner.unfoldDefinition}
\leanok
\uses{def:kernel-delta-value}
\texttt{isDelta} recognizes an expression whose head is a constant that has a delta value and is
applied to exactly \texttt{numLevelParams} universe arguments. \texttt{unfoldDefinitionCore}
replaces such a head constant by its value with the level parameters instantiated, memoizing the
result in the \texttt{unfold} cache when the level list is non-empty;
\texttt{unfoldDefinition} does the same at the head of an application and re-applies the
arguments.
\thesisref{axioms.tex \S2.5, $c_{\bar\ell} \rightsquigarrow v_{\bar\ell}(c)$}
\end{definition}

\begin{definition}[Primitive \texttt{Nat} reduction, \texttt{reduceNat}]
\label{def:kernel-reduce-nat}
\lean{Lean4Lean.TypeChecker.Inner.reduceNat, Lean4Lean.TypeChecker.Inner.reduceBinNatOp, Lean4Lean.TypeChecker.Inner.reduceBinNatPred, Lean4Lean.TypeChecker.Inner.reducePow, Lean4Lean.TypeChecker.Inner.rawNatLitExt?, Lean4Lean.TypeChecker.Inner.reduceNative, Lean4Lean.TypeChecker.Inner.reducePowMaxExp}
\leanok
The fast path for machine arithmetic: a unary \texttt{Nat.succ} application and the binary
\texttt{add}, \texttt{sub}, \texttt{mul}, \texttt{pow}, \texttt{gcd}, \texttt{mod}, \texttt{div},
\texttt{land}, \texttt{lor}, \texttt{xor}, \texttt{shiftLeft} and \texttt{shiftRight} reduce to a
\texttt{Nat} literal when their arguments normalize to literals, and \texttt{beq}/\texttt{ble} to
a \texttt{Bool} literal. \texttt{reducePow} refuses exponents above
$\texttt{reducePowMaxExp} = 2^{24}$. \texttt{reduceNative} throws on \texttt{reduceBool} and
\texttt{reduceNat} markers, which lean4lean does not support because that would require verified
compilation. \texttt{reduceNat} and \texttt{whnf} (\ref{def:kernel-whnf}) are mutually recursive;
the edge is recorded only in that other direction to keep the dependency graph acyclic.
\reviewnote{This is the one place where the checker's answers depend on the prelude really
defining \texttt{Nat.add} and friends correctly, which is why
Definition \ref{def:kernel-primitive-checkdef} exists. The exponent cap is a silent
incompleteness (a legitimate \texttt{Nat.pow} simply fails to reduce) and is not listed in
\texttt{divergences.md}, though the C++ kernel has the same cap.}
\end{definition}

\begin{definition}[Weak-head normalization without delta, \texttt{whnfCore'}]
\label{def:kernel-whnf-core}
\lean{Lean4Lean.TypeChecker.Inner.whnfCore', Lean4Lean.TypeChecker.Inner.whnfCore, Lean4Lean.TypeChecker.Inner.whnfFVar, Lean4Lean.TypeChecker.Inner.isLetFVar}
\leanok
\uses{def:kernel-reduce-recursor, def:kernel-reduce-proj}
Reduces to weak-head normal form using $\beta$ (iterated at the head, with the remaining arguments
re-applied), $\zeta$ (\texttt{let} and let-bound free variables, via \texttt{whnfFVar}), $\iota$
(through \texttt{reduceRecursor}) and projection reduction, but never unfolding definitions.
Results are memoized unless \texttt{cheapProj} is set, in which case the structure argument of a
projection is reduced only with \texttt{whnfCore} and nothing is cached.
\thesisref{axioms.tex \S2.3, the head reduction $e \rightsquigarrow^* k$; the $\beta$, $\zeta$ and $\iota$ rules}
\reviewnote{The docstring records that the C++ companion flag \texttt{cheap\_rec} is omitted
because nothing has set it since lean4\#9275, and the comment at
\srcloc{Lean4Lean/TypeChecker.lean}{404} justifies the re-decomposition branch. Both are the kind
of invariant note the proofs need.}
\end{definition}

\begin{definition}[Full weak-head normalization, \texttt{whnf'}]
\label{def:kernel-whnf}
\lean{Lean4Lean.TypeChecker.Inner.whnf', Lean4Lean.TypeChecker.Inner.whnf, Lean4Lean.TypeChecker.whnf}
\leanok
\uses{def:kernel-whnf-core, def:kernel-reduce-nat, def:kernel-delta-unfold}
Iterates \texttt{whnfCore'}, then \texttt{reduceNative} (which only ever throws or declines), then
\texttt{reduceNat}, then a single $\delta$-unfolding step, stopping when no rule applies. The loop
is bounded by \texttt{FuelConfig.whnf}, or \texttt{whnfEager} in an \texttt{eagerReduce} context,
and throws \texttt{.deterministicTimeout} on exhaustion. Easy syntactic cases are returned without
caching; everything else is memoized.
\thesisref{axioms.tex \S2.3, $e \rightsquigarrow^* k$; unique.tex \S4 (sec:kappa), $\rightsquigarrow_\kappa$ over $\beta\delta\zeta\iota$}
\end{definition}

\begin{definition}[Normalization to a sort or a pi type]
\label{def:kernel-ensure-sort-forall}
\lean{Lean4Lean.TypeChecker.Inner.ensureSortCore, Lean4Lean.TypeChecker.Inner.ensureForallCore}
\leanok
\uses{def:kernel-whnf}
\texttt{ensureSortCore e s} returns an $e' = \mathsf{Sort}\;\ell$ definitionally equal to $e$,
by weak-head normalization, and otherwise throws \texttt{typeExpected} blaming $s$;
\texttt{ensureForallCore} does the same for $\forall x:\alpha.\,\beta$, blaming the application
that required a function type.
\thesisref{axioms.tex \S2.1, the premises $\Gamma \vdash \alpha : \mathcal{U}_\ell$ of the $\forall$ and application rules}
\end{definition}

\section{Type inference}

\begin{definition}[Typing of constants and free variables]
\label{def:kernel-infer-constant}
\lean{Lean4Lean.TypeChecker.Inner.inferConstant, Lean4Lean.TypeChecker.Inner.checkLevel, Lean4Lean.TypeChecker.Inner.inferFVar}
\leanok
\uses{def:kernel-level-getundefparam, family:kernel-env-checks}
The type of $c_{\bar\ell}$ is the declaration's type instantiated at $\bar\ell$, after checking
that the number of levels matches \texttt{levelParams}. When \texttt{inferOnly} is false it
additionally rejects an unsafe constant used from non-unsafe code, a \texttt{partial} definition
used from safe code, and any $\ell_i$ mentioning a universe parameter outside $\bar u$.
\texttt{inferFVar} reads the type of a free variable out of $\Gamma$.
\thesisref{axioms.tex \S2.1 and \S2.5, $\vdash c_{\bar\ell} : \tau_{\bar\ell}(c)$ and $(x{:}\alpha) \in \Gamma$ implies $\Gamma \vdash x : \alpha$}
\reviewnote{The safety discipline \texttt{.safe}/\texttt{.partial}/\texttt{.unsafe} is outside the
thesis's system; it is the extra structure that the safety-indexed abstract environments of the
verification layer exist to model.}
\end{definition}

\begin{definition}[Typing of binders and applications]
\label{def:kernel-infer-binders}
\lean{Lean4Lean.TypeChecker.Inner.inferLambda, Lean4Lean.TypeChecker.Inner.inferForall, Lean4Lean.TypeChecker.Inner.inferLet, Lean4Lean.TypeChecker.Inner.inferApp}
\leanok
\uses{def:kernel-ensure-sort-forall, def:kernel-monad-localnamegen, def:kernel-cheap-beta}
Telescope-at-a-time typing. A lambda's type is the pi over its collected binders of the body's
type after \texttt{cheapBetaReduce}; a pi's type is $\mathcal{U}_{\mathrm{imax}(\ell_1,\dots,\ell_n,\ell)}$
folded with \texttt{mkLevelIMax'} over the binder levels; a \texttt{let} checks the value against
the ascribed type (when \texttt{inferOnly} is false) and binds a let-declaration;
\texttt{inferApp} walks an application spine in one pass, instantiating the pi bodies lazily with
\texttt{instantiateRevRange}.
\thesisref{axioms.tex \S2.1, the $\lambda$, $\forall$ and application rules; \S2.4 for \texttt{let}}
\reviewnote{\texttt{inferApp} is the \texttt{inferOnly := true} fast path and is correct only on
already well-typed terms; the precondition is recorded in the \texttt{isDefEqCore} docstring at
\srcloc{Lean4Lean/TypeChecker.lean}{172} rather than on \texttt{inferApp} itself.}
\end{definition}

\begin{definition}[Typing of structure projections, \texttt{inferProj}]
\label{def:kernel-infer-proj}
\lean{Lean4Lean.TypeChecker.Inner.inferProj}
\leanok
\uses{def:kernel-whnf}
The type of \texttt{.proj S i e}: normalize the type of $e$ to an application $I\;\bar a$, require
$I$ to be \texttt{S} and a single-constructor inductive with
$|\bar a| = \texttt{numParams} + \texttt{numIndices}$, instantiate the constructor's type at the
structure's levels and then at the parameters, walk $i$ binders instantiating each dependent field
with \texttt{.proj S j e}, and return the domain of the next binder. If the structure's universe
is not \texttt{isNeverZero}, every dependent field's domain and the result must additionally be a
proposition, checked with \texttt{isProp} (\ref{def:kernel-proof-irrel}). No thesis counterpart:
the thesis's type theory has no \texttt{proj}. This definition and \ref{def:kernel-proof-irrel}
are mutually recursive (through \ref{def:kernel-infer-type}); only the back-edges into this node
are recorded via \texttt{\textbackslash uses} to keep the dependency graph acyclic.
\reviewnote{A documented divergence: lean4lean tests \texttt{isNeverZero} where the C++ kernel
tests \texttt{!isAlwaysZero}, deliberately keeping \texttt{proj} no stronger than the recursor the
kernel generates for the same type. This directly bounds the scope of the \texttt{trproj}
\texttt{inferProj.WF} and \texttt{TrProj} results.}
\end{definition}

\begin{definition}[The algorithmic typing judgment, \texttt{inferType'}]
\label{def:kernel-infer-type}
\lean{Lean4Lean.TypeChecker.Inner.inferType', Lean4Lean.TypeChecker.Inner.inferType}
\leanok
\uses{def:kernel-infer-constant, def:kernel-infer-binders, def:kernel-infer-proj}
The main recursion. Loose bound variables are rejected outright and results are memoized in one of
two caches according to \texttt{inferOnly}. With \texttt{inferOnly := false} it implements Lean's
algorithmic typing judgment in full, checking level scopes, the existence of the literal types
\texttt{Nat}, \texttt{Char.ofNat} and \texttt{String.ofList}, and, at every application, that the
argument's type is definitionally equal to the domain (\ref{def:kernel-isdefeq-core}); with
\texttt{inferOnly := true} it assumes well-typedness and takes the \texttt{inferApp} fast path.
This definition and \ref{def:kernel-isdefeq-core} are mutually recursive; only the back-edge from
\ref{def:kernel-isdefeq-core} is recorded via \texttt{\textbackslash uses} to keep the dependency graph acyclic.
\thesisref{typesys.tex \S3, the algorithmic judgment $\Gamma \Vdash e : \alpha$ and the failure of subject reduction}
\reviewnote{The \texttt{eagerReduce} escape hatch at \srcloc{Lean4Lean/TypeChecker.lean}{299}
changes definitional-equality behaviour on a \emph{syntactic} test for core Lean's
\texttt{eagerReduce} marker in an argument, and is consulted again at lines 541, 782 and 845. It
mirrors upstream Lean, so it is rightly not a divergence, but it has no thesis counterpart and no
local documentation of what the four sites jointly guarantee.}
\end{definition}

\section{Definitional equality}

\begin{definition}[Reduction-free cases, \texttt{quickIsDefEq}]
\label{def:kernel-quick-isdefeq}
\lean{Lean4Lean.TypeChecker.Inner.quickIsDefEq, Lean4Lean.TypeChecker.Inner.isDefEqLambda, Lean4Lean.TypeChecker.Inner.isDefEqForall}
\leanok
\uses{def:kernel-equiv-manager, def:kernel-level-isequiv}
Decides the cases of $\Gamma \vdash t \Leftrightarrow s$ that need no reduction: $\alpha$-equivalence
or a previously recorded success through the \texttt{EquivManager}, two sorts by level
equivalence, two literals by equality, two \texttt{mdata} wrappers by their payloads, and two
lambdas or two pis by \texttt{isDefEqLambda}/\texttt{isDefEqForall}, which compare domains and
recurse under a delayed substitution. Everything else returns \texttt{.undef} and is deferred.
This definition and \texttt{isDefEqCore} (\ref{def:kernel-isdefeq-core}) are mutually recursive;
only the back-edge from \ref{def:kernel-isdefeq-core} is recorded via \texttt{\textbackslash uses} to keep the
dependency graph acyclic.
\thesisref{axioms.tex \S2.3, the congruence rules of $\Leftrightarrow$ and $\ell \equiv \ell'$ implies $\mathcal{U}_\ell \equiv \mathcal{U}_{\ell'}$}
\reviewnote{The \texttt{modifyGet} at \srcloc{Lean4Lean/TypeChecker.lean}{593} destructures
\texttt{TypeChecker.State} positionally as \texttt{.mk a1 \dots a7 (eqvManager := m)}; adding or
reordering a field silently changes what this matches. A named \texttt{modifyGet} would be
equivalent and robust.}
\end{definition}

\begin{definition}[Auxiliary definitional-equality rules]
\label{family:kernel-isdefeq-aux}
\lean{Lean4Lean.TypeChecker.Inner.isDefEqArgs, Lean4Lean.TypeChecker.Inner.isDefEqApp, Lean4Lean.TypeChecker.Inner.isDefEqOffset, Lean4Lean.TypeChecker.Inner.isDefEqUnitLike, Lean4Lean.TypeChecker.Inner.tryStringLitExpansionCore, Lean4Lean.TypeChecker.Inner.tryStringLitExpansion, Lean4Lean.TypeChecker.Inner.isNatZero, Lean4Lean.TypeChecker.Inner.isNatSuccOf?}
\leanok
\uses{def:kernel-lit-to-ctor}
The remaining special cases: congruence on application spines, with
\texttt{isDefEqApp} comparing argument counts before heads and \texttt{isDefEqArgs} assuming the
heads already agree and recursing through \texttt{isDefEqCore} (\ref{def:kernel-isdefeq-core});
the \texttt{Nat} offset rule reducing $n{+}1 \Leftrightarrow m{+}1$ to $n \Leftrightarrow m$
whether written as literals or \texttt{Nat.succ} applications; string-literal expansion against a
\texttt{String.ofList} application; and \texttt{isDefEqUnitLike}, which equates any two elements
of a non-recursive one-constructor type with no fields and no indices. This family and
\ref{def:kernel-isdefeq-core} are mutually recursive; only the back-edge from
\ref{def:kernel-isdefeq-core} is recorded via \texttt{\textbackslash uses} to keep the dependency graph acyclic.
\reviewnote{\texttt{isDefEqUnitLike} is not derivable from the thesis's rules; it is an extra
definitional equality the C++ kernel has, and the model has to account for it. Two entries of
\texttt{divergences.md} cover this family (argument-count-before-head, and the \texttt{whnf}
removed from \texttt{tryStringLitExpansionCore}).}
\end{definition}

\begin{definition}[Eta and structure-eta]
\label{def:kernel-eta-rules}
\lean{Lean4Lean.TypeChecker.Inner.tryEtaExpansionCore, Lean4Lean.TypeChecker.Inner.tryEtaExpansion, Lean4Lean.TypeChecker.Inner.tryEtaStructCore, Lean4Lean.TypeChecker.Inner.tryEtaStruct}
\leanok
\uses{def:kernel-infer-type, def:kernel-is-nonrec-structure}
Function $\eta$: when exactly one side is a lambda, compare it with
$\lambda x{:}A.\,s\;x$, where $(x:A) \to B$ is the inferred type of the other side. Structure
$\eta$: when one side is a saturated constructor application of a non-recursive structure and the
two sides have definitionally equal types, compare each constructor argument beyond the parameters
with the corresponding projection of the other side. Both are tried in each direction, recursing
into \texttt{isDefEqCore} (\ref{def:kernel-isdefeq-core}) on the compared subterms; this
definition and \ref{def:kernel-isdefeq-core} are mutually recursive, and only the back-edge from
\ref{def:kernel-isdefeq-core} is recorded via \texttt{\textbackslash uses} to keep the dependency graph acyclic.
\thesisref{axioms.tex \S2.2, the $\eta$ rule; structure $\eta$ has no thesis counterpart}
\reviewnote{The comment at \srcloc{Lean4Lean/TypeChecker.lean}{654} documents that both symmetric
calls in \texttt{tryEtaStruct} redo work \texttt{isDefEqApp} already did, and that the C++ kernel
has the same redundancy. Matching the kernel rather than optimizing is the right call here.}
\end{definition}

\begin{definition}[Proof irrelevance, \texttt{isDefEqProofIrrel}]
\label{def:kernel-proof-irrel}
\lean{Lean4Lean.TypeChecker.Inner.isDefEqProofIrrel, Lean4Lean.TypeChecker.Inner.isProp, Lean4Lean.TypeChecker.Inner.getSortLevel}
\leanok
\uses{def:kernel-infer-type, def:kernel-ensure-sort-forall}
Two terms are definitionally equal if the inferred type of the first is a proposition and the two
inferred types are definitionally equal (\ref{def:kernel-isdefeq-core}); otherwise the test
returns \texttt{.undef}. \texttt{getSortLevel} (\srcloc{Lean4Lean/TypeChecker.lean}{224}) infers a
type and normalizes it to a sort, and \texttt{isProp} asks whether that level
\texttt{isAlwaysZero}. This definition and \ref{def:kernel-isdefeq-core} are mutually recursive;
only the back-edge from \ref{def:kernel-isdefeq-core} is recorded via \texttt{\textbackslash uses} to keep the
dependency graph acyclic.
\thesisref{axioms.tex \S2.2, proof irrelevance; typesys.tex \S3 (sec:undecidable), why this makes $\equiv$ undecidable}
\reviewnote{\texttt{isProp} tests \texttt{isAlwaysZero} while \texttt{inferProj} and
\texttt{toCtorWhenStruct} test \texttt{isNeverZero} for the opposite question. The asymmetry is
deliberate and explained in \texttt{divergences.md}, but the two spellings sitting in one file
invite confusion.}
\end{definition}

\begin{definition}[Lazy delta reduction, \texttt{lazyDeltaReduction}]
\label{def:kernel-lazy-delta}
\lean{Lean4Lean.TypeChecker.Inner.lazyDeltaReductionStep, Lean4Lean.TypeChecker.Inner.lazyDeltaReduction, Lean4Lean.TypeChecker.Inner.lazyDeltaProjReduction, Lean4Lean.TypeChecker.Inner.tryUnfoldProjApp, Lean4Lean.TypeChecker.Inner.failedBefore, Lean4Lean.TypeChecker.Inner.cacheFailure, Lean4Lean.TypeChecker.ReductionStatus, Lean4Lean.TypeChecker.ReductionStatus.bool}
\leanok
\uses{def:kernel-delta-unfold, def:kernel-quick-isdefeq, family:kernel-isdefeq-aux, axiom:kernel-ptr-eq, def:kernel-reducibility-lt, def:kernel-reduce-nat}
The heuristic heart of definitional equality. Given two \texttt{cheapProj} weak-head normal forms,
each step tries the \texttt{Nat} offset rule, then literal and native reduction, then one
\texttt{lazyDeltaReductionStep}: unfold whichever side has the strictly more reducible head by
\texttt{ReducibilityHints.lt'}, retrying an unreduced projection head with \texttt{tryUnfoldProjApp}
when only one side is $\delta$-reducible; and when both heads are the same regular definition (a
test that uses \texttt{ptrEqConstantInfo}), first compare level arguments and then arguments
congruentially, caching the failure before unfolding both. The loop is bounded by
\texttt{FuelConfig.lazyDelta}. \texttt{lazyDeltaProjReduction} is the analogue for two projections
at the same index, falling back to comparing the projected fields.
\thesisref{axioms.tex \S2.5, $\delta$ inside $\Leftrightarrow$; typesys.tex \S3, algorithmic equality is not transitive}
\reviewnote{\texttt{ptrEqConstantInfo} is used here as a \emph{semantic} test, not an
optimization: two different addresses for the same \texttt{ConstantInfo} select a different
branch. That is precisely what the \texttt{PtrEq} docstring warns about and why
\texttt{ptrEqConstantInfo\_eq} must be an axiom.}
\end{definition}

\begin{definition}[Algorithmic definitional equality, \texttt{isDefEqCore'}]
\label{def:kernel-isdefeq-core}
\lean{Lean4Lean.TypeChecker.Inner.isDefEqCore', Lean4Lean.TypeChecker.Inner.isDefEqCore, Lean4Lean.TypeChecker.Inner.isDefEq, Lean4Lean.TypeChecker.isDefEq}
\leanok
\uses{def:kernel-quick-isdefeq, def:kernel-whnf-core, def:kernel-proof-irrel, def:kernel-lazy-delta, family:kernel-isdefeq-aux, def:kernel-eta-rules, axiom:kernel-ptr-eq}
The decision procedure for $\Gamma \vdash t \Leftrightarrow s$, in order: the quick cases with
hashing; the asymmetric shortcut normalizing $t$ when $s$ is the constant \texttt{true};
\texttt{cheapProj} weak-head normalization of both sides followed by the quick cases again (skipped
when \texttt{ptrEqExpr} says neither side moved); proof irrelevance; lazy delta reduction;
syntactic equality of constants with equivalent levels, of free variables, and of projections at
the same index (which hands over to \texttt{lazyDeltaProjReduction}); a full \texttt{whnfCore} of
both sides with a restart if a projection reduced; and finally congruence on applications,
$\eta$, structure $\eta$, string-literal expansion and the unit-like rule.
\thesisref{axioms.tex \S2.3, $\Gamma \vdash e \Leftrightarrow e'$; typesys.tex \S3, item:alg\_defn}
\reviewnote{This is the function whose correctness is the entire point of the verification. The
\texttt{s.isConstOf ``true} shortcut at \srcloc{Lean4Lean/TypeChecker.lean}{845} is an asymmetric
special case for \texttt{Decidable} instances with neither a thesis counterpart nor a
\texttt{divergences.md} entry.}
\end{definition}

\section{Type checker entry points}

\begin{definition}[Type checker entry points]
\label{family:kernel-tc-entrypoints}
\lean{Lean4Lean.TypeChecker.whnfCore, Lean4Lean.TypeChecker.unfoldDefinition, Lean4Lean.TypeChecker.inferType, Lean4Lean.TypeChecker.checkType, Lean4Lean.TypeChecker.isProp, Lean4Lean.TypeChecker.ensureSort, Lean4Lean.TypeChecker.ensureForall, Lean4Lean.TypeChecker.ensureType, Lean4Lean.TypeChecker.M.runTermElab, Lean4Lean.TypeChecker.RecM.runTermElab}
\leanok
\uses{structure:kernel-methods-fuel}
The \texttt{M}-level wrappers that run the fuel-bounded \texttt{RecM} operations at the configured
recursion depth, together with \texttt{checkType} (inference with \texttt{inferOnly := false}),
\texttt{ensureType}, and the \texttt{TermElabM} lifts used to call the kernel from inside Lean
elaboration for testing.
\reviewnote{A commented-out \texttt{example}/\texttt{run\_tac} block at
\srcloc{Lean4Lean/TypeChecker.lean}{968} is leftover scratch code that should be either a real
test or deleted.}
\end{definition}

\begin{definition}[Eta expansion to full arity, \texttt{etaExpand}]
\label{def:kernel-eta-expand}
\lean{Lean4Lean.TypeChecker.etaExpand}
\leanok
\uses{def:kernel-whnf, family:kernel-tc-entrypoints, def:kernel-monad-localnamegen}
Walks under a term's existing binders, then repeatedly normalizes its type and adds one binder per
pi, bounded by \texttt{FuelConfig.etaExpand}, finally rebuilding the lambda over all collected
variables applied to the accumulated arguments. If the type is not a pi after the first
normalization the original term is returned.
\thesisref{unique.tex \S4 (sec:kappa), the $\eta$-expansion map $\bar e$ and $\mathrm{rec}$-normal form}
\reviewnote{This function has no call site anywhere under \texttt{Lean4Lean/}. If it is meant to
realize the thesis's $\mathrm{rec}$-normal-form preprocessing, which saturates every
$\mathrm{rec}$ and $\mathrm{lift}$ before $\kappa$-reduction is even defined, that intent is
undocumented, and the executable kernel never enforces $\mathrm{rec}$-normal form. This matters:
the saturation hypothesis of \texttt{inductiveReduceRecCore.WF} is the executable counterpart of
exactly that preprocessing.}
\end{definition}

\section{Inductive declarations}

\begin{definition}[State of the inductive-block adder]
\label{structure:kernel-addinductive-ctx}
\lean{Lean4Lean.AddInductive.RecInfo, Lean4Lean.AddInductive.InductiveStats, Lean4Lean.AddInductive.Context, Lean4Lean.AddInductive.M}
\leanok
\uses{structure:kernel-tc-monad, def:kernel-monad-localnamegen}
\texttt{InductiveStats} records what the block-wide checks share: the local context, the universe
arguments and the result level $\ell$, the per-type index counts, the constants standing for the
block's types, the free variables for the parameters $a{::}\alpha$, and whether $\ell$ is
provably nonzero. \texttt{RecInfo} records one recursor's motive, minor premises, indices and
major premise. \texttt{M} is \texttt{ReaderT Context (Except Exception)} with a \texttt{MonadLift}
from \texttt{TypeChecker.M}, so a type checker action runs in the block's context.
\thesisref{axioms.tex \S2.6 (sec:inductive), the labelling $F = \forall a{::}\alpha.\ \mathcal{U}_\ell$ and $K = \sum_c (c : \forall b{::}\beta.\ t\;p[b])$}
\end{definition}

\begin{definition}[Checking and declaring the block's type constants]
\label{def:kernel-check-inductive-types}
\lean{Lean4Lean.AddInductive.checkInductiveTypes, Lean4Lean.AddInductive.declareInductiveTypes, Lean4Lean.AddInductive.isRec, Lean4Lean.AddInductive.isReflexive, Lean4Lean.AddInductive.hasIndOcc}
\leanok
\uses{family:kernel-tc-entrypoints, structure:kernel-addinductive-ctx}
Type-checks each declared family, walks its telescope requiring the first \texttt{nparams} binders
to be definitionally equal to the parameters fixed by the first type, counts the remaining binders
as indices, and requires the result to be a sort whose level is equivalent across the whole block.
It then declares the block's type constants in the environment so that the constructors can be
checked against them. \texttt{isRec} and \texttt{isReflexive} compute the flags stored in the
\texttt{inductInfo}, both built on \texttt{hasIndOcc}.
\thesisref{axioms.tex \S2.6.1 (sec:inductive), the $\mathsf{spec}$ judgment}
\end{definition}

\begin{definition}[Strict positivity check, \texttt{checkPositivity}]
\label{def:kernel-check-positivity}
\lean{Lean4Lean.AddInductive.checkPositivity, Lean4Lean.AddInductive.isRecArg, Lean4Lean.AddInductive.isValidIndApp?, Lean4Lean.AddInductive.isValidIndAppIdx}
\leanok
\uses{def:kernel-check-inductive-types}
Walks a constructor argument type: if it does not mention the block at all it is accepted; if it
is a pi whose domain mentions the block it is rejected as a non-positive occurrence; otherwise the
argument must end in a valid application $t\;\bar e$ of one of the block's types to exactly the
block's parameters, with no further occurrence of the block among the indices
(\texttt{isValidIndAppIdx}). The loop is bounded by \texttt{FuelConfig.inductiveFuel}.
\thesisref{axioms.tex \S2.6.1 (sec:inductive), the two $\mathsf{ctor}$ rules}
\reviewnote{The model-side counterpart on the \texttt{iota} branch is the positivity
specification in \texttt{Theory/Inductive.lean} (Chapter \ref{chap:inductive}); the two should be
cross-checked on nested and reflexive cases, which the kernel handles in
\texttt{ElimNestedInductive} rather than in this judgment.}
\end{definition}

\begin{definition}[Checking and declaring constructors]
\label{def:kernel-check-constructors}
\lean{Lean4Lean.AddInductive.checkConstructors, Lean4Lean.AddInductive.declareConstructors}
\leanok
\uses{def:kernel-check-positivity, family:kernel-tc-entrypoints}
For each constructor: reject a duplicate name, type-check its type, require its first binders to
match the block's parameters, require every later argument's universe $\ell'$ to satisfy
$\ell' \le \ell$ unless $\ell$ is always zero (using the complete \texttt{geq'}), run the
positivity check on it unless the block is unsafe, and require the result to be a valid
application of the constructor's own type. \texttt{declareConstructors} then adds each
\texttt{ctorInfo} with its parameter count and field count $= \text{arity} - \texttt{numParams}$.
\thesisref{axioms.tex \S2.6.1 (sec:inductive), the $\mathsf{ctor}$ judgment and its universe side condition}
\end{definition}

\begin{definition}[Large elimination and K-likeness]
\label{def:kernel-large-elim}
\lean{Lean4Lean.AddInductive.isLargeEliminator, Lean4Lean.AddInductive.getElimLevel, Lean4Lean.AddInductive.isKTarget}
\leanok
\uses{family:kernel-tc-entrypoints}
\texttt{isLargeEliminator} is true when the result universe is provably nonzero, or when the block
is a single type with no constructors, or with one constructor all of whose non-parameter
arguments are either propositions or appear among the arguments of the output type (the thesis's
$y \in e$ side condition). \texttt{getElimLevel} then returns a fresh universe variable, and
$0$ otherwise. \texttt{isKTarget} is true for a single type in an always-zero universe with one
constructor that has no fields beyond the parameters.
\thesisref{axioms.tex \S2.6.2 (sec:large\_elim), the $\mathsf{LE}$ and $\mathsf{LE}$-$\mathsf{ctor}$ judgments}
\end{definition}

\begin{definition}[Motives and minor premises, \texttt{mkRecInfos}]
\label{def:kernel-mk-rec-infos}
\lean{Lean4Lean.AddInductive.mkRecInfos, Lean4Lean.AddInductive.mkRecInfos.loopArgs1, Lean4Lean.AddInductive.mkRecInfos.loopInd1, Lean4Lean.AddInductive.mkRecInfos.loopCtorArgs, Lean4Lean.AddInductive.mkRecInfos.loopUArgs, Lean4Lean.AddInductive.mkRecInfos.loopU, Lean4Lean.AddInductive.mkRecInfos.loopCtors, Lean4Lean.AddInductive.mkRecInfos.loopInd2}
\leanok
\uses{def:kernel-large-elim, def:kernel-monad-localnamegen}
Builds, per type of the block, the indices, the major premise $t : I\;\bar a\;\bar i$ and the
motive $C : \forall \bar i,\ I\;\bar a\;\bar i \to \mathcal{U}_u$ at the elimination level, and
then, per constructor, the minor premise
$\varepsilon_c = \forall b{::}\beta,\ \forall v{::}\delta,\ C\;p[b]\;(c\;b)$ with
$\delta_i = \forall x{::}\xi_i,\ C\;\pi_i[b,x]\;(u_i\;x)$ for each recursive argument, where the
recursive arguments are those \texttt{isRecArg} accepts.
\thesisref{axioms.tex \S2.6.3, the definitions of $\kappa$, $\varepsilon$ and $\delta$}
\reviewnote{A tower of continuation-passing loops with hand-written \texttt{termination\_by}, split
into standalone functions because they would not compile together as \texttt{let rec}s. It is
readable only against the thesis's notation, and there is no docstring pointing at \S2.6.3.}
\end{definition}

\begin{definition}[Generation of the $\iota$ rules, \texttt{mkRecRules}]
\label{def:kernel-mk-rec-rules}
\lean{Lean4Lean.AddInductive.mkRecRules, Lean4Lean.AddInductive.getRecLevels, Lean4Lean.AddInductive.getRecLevelParams}
\leanok
\uses{def:kernel-mk-rec-infos}
For each constructor, builds the \texttt{RecursorRule} with \texttt{nfields} the number of
non-parameter constructor arguments and right-hand side
$\lambda \bar a\,\bar C\,\bar e\,b{::}\beta.\ e_c\;b\;v$ where
$v_i = \lambda x{::}\xi_i.\ \mathrm{rec}_{P_j}\;\bar a\;\bar C\;\bar e\;\pi_i[b,x]\;(u_i\;x)$, at
the recursor's levels (the elimination level prepended when it is a parameter). These rules are
exactly what \texttt{inductiveReduceRecCore} later consumes.
\thesisref{axioms.tex \S2.6.4 (sec:iota), $\mathrm{rec}_P\;C\;e\;p[b]\;(c\;b) \equiv e_c\;b\;v$}
\reviewnote{This is the producer side of the contract whose consumer side
\texttt{inductiveReduceRecCore.WF} verifies, and the \texttt{iota} branch's model of the $\iota$
rule is a model of exactly these rules; it is the right place to look when judging whether that
specification is faithful.}
\end{definition}

\begin{definition}[Adding a non-nested inductive block, \texttt{AddInductive.run}]
\label{def:kernel-addinductive-run}
\lean{Lean4Lean.AddInductive.run}
\leanok
\uses{def:kernel-check-inductive-types, def:kernel-check-constructors, def:kernel-large-elim, def:kernel-mk-rec-infos, def:kernel-mk-rec-rules}
The whole pipeline for one block: reject duplicate universe parameters, check and declare the
types, check and declare the constructors, compute the elimination level and the K-like flag,
build the motives and minor premises, and add one \texttt{recInfo} per type whose type is
$\forall \bar a,\ \bar C,\ \bar e,\ \bar i,\ (z : I\;\bar a\;\bar i),\ C\;\bar i\;z$ and whose
rules come from \texttt{mkRecRules}.
\thesisref{axioms.tex \S2.6 (sec:inductive) in its entirety}
\reviewnote{The binding \texttt{isUnsafe}, read from the context safety, is computed twice, at
\srcloc{Lean4Lean/Inductive/Add.lean}{453} and again at
\srcloc{Lean4Lean/Inductive/Add.lean}{470}, the second shadowing the first. The recursor type is
built with \texttt{type.inferImplicit 1000 false}, a magic constant whose accompanying comment
notes that the flag's polarity is inverted relative to C++.}
\end{definition}

\begin{definition}[Nested inductive elimination, \texttt{ElimNestedInductive.run}]
\label{def:kernel-elim-nested}
\lean{Lean4Lean.ElimNestedInductive.run, Lean4Lean.ElimNestedInductive.Result, Lean4Lean.ElimNestedInductive.Result.restoreNested, Lean4Lean.ElimNestedInductive.Result.restoreCtorName, Lean4Lean.ElimNestedInductive.replaceIfNested, Lean4Lean.ElimNestedInductive.replaceAllNested, Lean4Lean.ElimNestedInductive.isNestedInductiveApp?, Lean4Lean.ElimNestedInductive.mkUniqueName}
\leanok
Rewrites a nested inductive block into a mutual one: each nested occurrence inside a constructor
type is replaced by a fresh auxiliary inductive declared in the same block, with the correspondence
recorded in \texttt{aux2nested}; \texttt{restoreNested} and \texttt{restoreCtorName} invert the
rewrite afterwards on the generated constructor types, recursor types and rule right-hand sides.
No thesis counterpart: the thesis's $\mu\,t{:}F.\,K$ has no nested inductives.
\reviewnote{Fully unverified and the largest source of \texttt{unreachable!}/\texttt{assert!} in
the kernel. Three \texttt{divergences.md} entries concern it, one of which admits an actual hole:
\texttt{restoreCtorName}'s \texttt{unreachable!} would continue with a default value, and a
restored rule constructor \emph{name} is covered by no later re-check, since lean4lean also skips
lean4\#14621's recheck of restored declarations.}
\end{definition}

\begin{definition}[The inductive declaration entry point, \texttt{addInductive}]
\label{def:kernel-add-inductive}
\lean{Lean4Lean.Environment.addInductive, Lean4Lean.mkAuxRecNameMap, Lean4Lean.checkNoNestedAux}
\leanok
\uses{def:kernel-elim-nested, def:kernel-addinductive-run, family:kernel-tc-entrypoints}
Checks the block's types and constructor types for metavariables and free variables and its
constructor types for the reserved \texttt{\_nested} prefix, runs nested-inductive elimination,
adds the resulting block with \texttt{AddInductive.run}, and, when nesting occurred, restores the
nested occurrences in the constructor and recursor declarations, renames the auxiliary recursors
through \texttt{mkAuxRecNameMap}, and finally type-checks the recorded nested applications.
\thesisref{axioms.tex \S2.6 (sec:inductive)}
\reviewnote{This is the function for which \texttt{addDecl.WF}'s \texttt{inductDecl} case is
\texttt{sorry} at \srcloc{Lean4Lean/Verify/Environment.lean}{208}: the \texttt{AddInduct} witness
the model needs has to be constructed from here, and nobody has. Everything the \texttt{iota}
branch proves about inductive blocks sits on the model side of that gap.}
\end{definition}

\section{Primitive declarations}

\begin{definition}[The primitive recognizer's vocabulary]
\label{family:kernel-primitive-dsl}
\lean{Lean4Lean.Primitive.Reflection, Lean4Lean.Primitive.Condition, Lean4Lean.Primitive.ConditionImpl, Lean4Lean.Primitive.Probe, Lean4Lean.Primitive.unfoldWellFounded, Lean4Lean.Primitive.unfoldNatWellFounded, Lean4Lean.Primitive.checkNatFuelRec, Lean4Lean.Primitive.checkNatBoolCases, Lean4Lean.Primitive.microQq, Lean4Lean.lambdaTelescope}
\leanok
\uses{def:kernel-isdefeq-core, def:kernel-whnf}
The small language the primitive checks are written in: \texttt{Reflection} describes the two
shapes of \texttt{Decidable}-style reflection the prelude may use, \texttt{Condition} packages a
decidable \texttt{Nat} predicate with its \texttt{ite}/\texttt{dite}/\texttt{decide} forms,
\texttt{Probe} and \texttt{unfoldWellFounded} certify that a definition is the compiled form of a
well-founded recursion and expose its equation body, and the \texttt{q(\dots)} elaborator (a
hand-rolled mini-Qq built on \texttt{deriving ToExpr for Expr}) writes the expected expressions.
\reviewnote{\texttt{deriving instance ToExpr} for \texttt{Syntax}, \texttt{KVMap} and \texttt{Expr}
ties this file to the exact shape of core's \texttt{Expr} and \texttt{Syntax}. Acceptable for a
recognizer, but it is compile-time-only machinery living inside the trusted kernel library.}
\end{definition}

\begin{definition}[Primitive definition recognizer, \texttt{Primitive.checkDef}]
\label{def:kernel-primitive-checkdef}
\lean{Lean4Lean.Primitive.checkDef}
\leanok
\uses{family:kernel-primitive-dsl, def:kernel-isdefeq-core}
Dispatches on the declaration's name: for each of \texttt{Nat.add}, \texttt{pred}, \texttt{sub},
\texttt{mul}, \texttt{pow}, \texttt{mod}, \texttt{div}, \texttt{gcd}, \texttt{beq}, \texttt{ble},
\texttt{bitwise}, \texttt{land}, \texttt{lor}, \texttt{xor}, \texttt{shiftLeft},
\texttt{shiftRight}, and for \texttt{Char.ofNat} and \texttt{String.ofList}, it checks that the
definition actually supplied has the type and the definitional behaviour the GMP fast paths
assume. It returns \texttt{true} when the name is primitive (which is what permits the reserved
name) and \texttt{false} otherwise, and throws when the shape is wrong. No thesis counterpart:
the thesis's system has no primitive arithmetic.
\reviewnote{A genuine soundness gap in the C++ kernel that lean4lean closes, and the
\texttt{run\_meta} self-test at the end of the file gives real assurance that it accepts the
actual prelude. But the property it establishes is stated nowhere here; the connection to
\texttt{reduceNat} is made only in \texttt{Verify/Environment/Primitive/}
(Chapters \ref{chap:primitives-core} and \ref{chap:primitives-arith}), and a build-time test is
not a proof.}
\end{definition}

\begin{definition}[Primitive inductive recognizer, \texttt{Primitive.checkInductive}]
\label{def:kernel-primitive-checkinductive}
\lean{Lean4Lean.Primitive.checkInductive}
\leanok
Returns \texttt{true} only for a safe, parameterless, non-universe-polymorphic single inductive in
\texttt{Type} named \texttt{Bool} with constructors \texttt{false} then \texttt{true}, or
\texttt{Nat} with \texttt{zero} then \texttt{succ : Nat} $\to$ \texttt{Nat}; anything else with
those names throws. Constructor \emph{order} is checked, which matters because the literal
representation and the constructor-order-sensitive fast paths depend on it, and because
\texttt{Bool}'s constructor order is a known cross-system trap.
\end{definition}

\section{Declaration admission}

\begin{definition}[Per-kind declaration checkers]
\label{family:kernel-add-decls}
\lean{Lean4Lean.checkConstantValBody, Lean4Lean.checkConstantVal, Lean4Lean.checkDefinitionBody, Lean4Lean.addAxiom, Lean4Lean.addDefinition, Lean4Lean.addTheorem, Lean4Lean.addOpaque, Lean4Lean.addMutual}
\leanok
\uses{family:kernel-tc-entrypoints, def:kernel-isdefeq-core, def:kernel-primitive-checkdef, family:kernel-env-checks}
Each kind is checked before being committed. An axiom needs its type to be a well-formed type; a
definition additionally needs its body's type to be definitionally equal to the ascribed one, and
an \texttt{unsafe} definition is checked in an environment where the constant has been added as an
axiom of the same type (so its body may refer to it but not unfold it); a theorem additionally
needs its type to be a proposition; an opaque is like a definition; a mutual block, which must be
unsafe or partial, is checked with all its headers present as axioms, after requiring matching
safety, matching universe parameters and distinct names.
\thesisref{axioms.tex \S2.5, $\mathsf{constant}\;c_{\bar u}:\alpha$ and $\mathsf{def}\;c_{\bar u}:\alpha := e$}
\reviewnote{Three \texttt{divergences.md} entries live here. The comment at
\srcloc{Lean4Lean/Environment.lean}{53} explains why \texttt{Primitive.checkDef} runs after the
body check (so that its \texttt{isDefEq} calls see terms already known well-typed); that is an
ordering invariant the proofs depend on and it is easy to break.}
\end{definition}

\begin{definition}[The kernel's entry point, \texttt{addDecl}]
\label{def:kernel-add-decl}
\lean{Lean4Lean.addDecl}
\leanok
\uses{family:kernel-add-decls, def:kernel-add-quot, def:kernel-add-inductive, def:kernel-primitive-checkinductive}
Dispatches a \texttt{Declaration} to the checker for its kind, returning the extended environment
or a \texttt{Kernel.Exception}. For an inductive declaration it first consults
\texttt{Primitive.checkInductive} to decide whether reserved names are permitted. This is the
single function the whole verification effort is about: \texttt{addDecl.WF} states that a
successful call preserves well-formedness and extends every safety-indexed abstract environment.
\thesisref{soundness.tex \S6 (thm:sound) is the thesis's corresponding statement, and it has no counterpart in this repository}
\reviewnote{\texttt{addDecl.WF} is proved for every case except \texttt{inductDecl}, which is
\texttt{sorry} at \srcloc{Lean4Lean/Verify/Environment.lean}{208}. The \texttt{quotDecl} case is
unconditional thanks to the \texttt{iota} change to \texttt{checkEqType}.}
\end{definition}

\section{Replay and the command line driver}

\begin{definition}[Constant-by-constant replay]
\label{family:kernel-replay-core}
\lean{Lean4Lean.Replay.Context, Lean4Lean.Replay.State, Lean4Lean.Replay.M, Lean4Lean.Replay.isTodo, Lean4Lean.Replay.addDecl, Lean4Lean.Replay.replayConstant, Lean4Lean.Replay.replayConstants, Lean4Lean.Replay.throwKernelException}
\leanok
\uses{def:kernel-add-decl}
The driver that re-adds a module's constants to a kernel environment in dependency order:
\texttt{replayConstant} first replays the constants a declaration mentions, reconstructs a
\texttt{Declaration} from the stored \texttt{ConstantInfo}, and sends it to
\texttt{Lean4Lean.addDecl}; constructors and recursors are postponed rather than re-added, since
\texttt{addInductive} regenerates them. Unverified glue by design.
\reviewnote{The \texttt{--compare} path times lean4lean against Lean's own kernel, which is useful
but makes the driver depend on core's \texttt{Environment.addDecl} as well.}
\end{definition}

\begin{definition}[Post-replay consistency checks]
\label{family:kernel-replay-postponed}
\lean{Lean4Lean.Replay.checkPostponedConstructors, Lean4Lean.Replay.checkPostponedRecursors, Lean4Lean.Replay.checkQuotInit}
\leanok
\uses{family:kernel-replay-core}
After all constants are replayed, every postponed constructor and recursor that
\texttt{addInductive} regenerated is compared for \texttt{BEq} equality with the one stored in the
\texttt{.olean}, and (unless disabled) the environment is required to have \texttt{quotInit} set by
the end of the module, which is the assumption \texttt{finalizeImport} builds in. A mismatch is
reported as a plain \texttt{IO.userError}.
\reviewnote{This is the real value of replay as a test: a differential check of
\texttt{AddInductive.mkRecRules} against Lean's own elaborator across whole libraries, and in
practice the only end-to-end evidence that the kernel's rule generation agrees with Lean's. It is
worth pointing at when judging the \texttt{iota} branch's $\iota$-rule specification.}
\end{definition}

\begin{definition}[Replay entry points]
\label{def:kernel-replay-entrypoints}
\lean{Lean4Lean.Replay.replay, Lean4Lean.Replay.replayFromImports, Lean4Lean.Replay.replayFromFresh}
\leanok
\uses{family:kernel-replay-core, family:kernel-replay-postponed, def:kernel-env-empty-import}
\texttt{replay} runs the driver over a set of new constants, skipping unsafe and partial ones;
\texttt{replayFromImports} reconstructs the environment as it stood at the start of a module (via
\texttt{finalizeImport}) and replays only that module's own constants, freeing the compacted
regions afterwards under a careful use-after-free argument;
\texttt{replayFromFresh} replays a module and everything it imports into a completely empty
environment.
\reviewnote{\texttt{replayFromFresh} passes \texttt{checkQuot := false}
(\srcloc{Lean4Lean/Replay.lean}{319}), so the quotient-initialization check is skipped exactly in
the mode that builds an environment from scratch. In that mode \texttt{Init.Prelude} is replayed
and so \texttt{addQuot} does run, but the ordering is the opposite of what one would want.}
\end{definition}

\begin{definition}[Fuel configuration from the command line]
\label{family:kernel-fuelconfig-cli}
\lean{Lean4Lean.FuelConfig.ofJson?, Lean4Lean.FuelConfig.ofFile, Lean4Lean.FuelConfig.applyFlag, Lean4Lean.FuelConfig.toObj, Lean4Lean.FuelConfig.fieldNames}
\leanok
\uses{structure:kernel-fuel-config}
Layering of a JSON object, read from a file or built from a single \texttt{--config:field=value}
override, over an existing \texttt{FuelConfig}: unknown fields are rejected up front, known ones
are overlaid on the base config's own serialization, and the merged object goes back through the
derived \texttt{fromJson?} so that value validation stays in one place. The accepted field names
(\texttt{fieldNames}) are derived from the record's own encoding, so the two cannot drift apart.
\reviewnote{\texttt{toObj} calls \texttt{panic!} if the derived encoding is not a JSON object.
Unreachable, but a panic inside a command line argument parser. The node map named these two
helpers \texttt{toJsonObj} and \texttt{fields}; the declarations are actually called
\texttt{toObj} and \texttt{fieldNames}, and both are \texttt{private} to \texttt{Main}.}
\end{definition}

\begin{definition}[The lean4lean executable, \texttt{main}]
\label{def:kernel-main}
\lean{main, getCurrentModule}
\leanok
\uses{def:kernel-replay-entrypoints, family:kernel-fuelconfig-cli}
Parses the flags \texttt{-v}, \texttt{--fresh}, \texttt{--compare}, \texttt{--config=FILE} and
\texttt{--config:field=value}, resolves the target module names against the search path (defaulting
to the current Lake project's main module, read from the manifest by \texttt{getCurrentModule}),
and replays each target, in parallel \texttt{IO.asTask}s except under \texttt{--fresh}, reporting
the number of declarations checked. Unverified glue.
\end{definition}

\section{Contribution summary}

The two branches together changed 22 of the 4040 lines in this chapter's scope, and both changes
were made to unblock a specific proof obligation rather than to improve the kernel for its own
sake.

On \texttt{trproj}, the $\iota$ step of \texttt{inductiveReduceRec} was extracted verbatim into the
pure function \texttt{inductiveReduceRecCore} (Definition \ref{def:kernel-iota-core}), and
\texttt{inductiveReduceRec} (Definition \ref{def:kernel-inductive-reduce-rec}) now ends by calling
it. The reason is that \texttt{inductiveReduceRec} is monadic only because of its K-like
(Definition \ref{def:kernel-to-ctor-when-k}) and structure-$\eta$
(Definition \ref{def:kernel-to-ctor-when-struct}) preludes, whereas the $\iota$ step itself is
pure and total; after the split, \texttt{inductiveReduceRecCore.WF} can be stated about a function
with no monad, no environment lookup and no well-typedness side conditions beyond its hypotheses.
The commit is explicit that \texttt{reduceRecursor.WF}
(Definition \ref{def:kernel-reduce-recursor}) stays open, because the two preludes still have no
\texttt{IsDefEq} counterpart in the model. The split is behaviour-preserving; the extracted body
was diffed against master line by line and differs only in recomputing \texttt{rval.getMajorIdx}
where the caller had bound it.

On \texttt{iota}, \texttt{checkEqType} (Definition \ref{def:kernel-check-eq-type}) gained one line
rejecting an \texttt{unsafe} \texttt{Eq}. This is a reject-more divergence from \texttt{quot.cpp},
recorded in \texttt{divergences.md}, and it discharges rather than adds an assumption: both
\texttt{addQuot.WF} and \texttt{addDecl.WF} (Definition \ref{def:kernel-add-decl}) lost their
\texttt{Environment.EqSafe} hypothesis, and the assumption itself was deleted.

What remains open on the kernel side is unchanged by either branch and bounds what they can claim:
\texttt{addDecl.WF}'s \texttt{inductDecl} case is still \texttt{sorry} because the model witness
must be built from \texttt{Environment.addInductive}
(Definition \ref{def:kernel-add-inductive}); \texttt{reduceRecursor.WF} is open; the projection
rules (Definitions \ref{def:kernel-infer-proj} and \ref{def:kernel-reduce-proj}) carry a documented
\texttt{isNeverZero} divergence that narrows every \texttt{TrProj} statement; and
\texttt{etaExpand} (Definition \ref{def:kernel-eta-expand}), the only plausible executable
counterpart of the thesis's $\mathrm{rec}$-normal-form preprocessing, is never called, so the
saturation hypothesis of \texttt{inductiveReduceRecCore.WF} is not enforced anywhere in the kernel.

\section{Review notes}

\begin{itemize}
\item \textbf{Medium.} \srcloc{Lean4Lean/Quot.lean}{24}: the \texttt{iota} safety check is well
justified and net positive, but the guarantee is narrower than the commit message suggests.
\texttt{Environment.addQuot} returns early when \texttt{env.quotInit} is already true, and
\texttt{Kernel.Environment.finalizeImport} (\srcloc{Lean4Lean/Environment/Basic.lean}{124}) sets
\texttt{quotInit := !imports.isEmpty} unconditionally, so on any imported environment neither the
shape check nor the new safety check ever runs. The model-side claim holds only for environments
replayed from scratch.
\item \textbf{Low.} \srcloc{Lean4Lean/Inductive/Reduce.lean}{71}:
\texttt{inductiveReduceRecCore} is declared inside the \texttt{section} opened at line 9 with
variables \texttt{[Monad m]}, \texttt{env}, \texttt{whnf}, \texttt{inferType},
\texttt{isDefEq}, none of which it uses. Lean binds only used section variables, so the signature
is the intended four-argument one, but a reader (or a later edit that adds a use of \texttt{env})
can be misled about its arity. Placing it outside the section would be clearer.
\item \textbf{Low.} \srcloc{Lean4Lean/Inductive/Reduce.lean}{91}: after the split, the unchanged
master docstring of \texttt{inductiveReduceRec} still describes the $\iota$ slicing in full, which
is now \texttt{inductiveReduceRecCore}'s job and is documented there too. The file states the same
algorithm twice; the caller's docstring should have been trimmed to the K-like, structure-$\eta$
and literal preprocessing it still performs.
\item \textbf{Low.} \srcloc{Lean4Lean/Inductive/Reduce.lean}{75}:
\texttt{inductiveReduceRecCore} collapses ``no rule for this constructor'', ``major has too few
arguments'' and ``wrong number of universe arguments'' into a single \texttt{none}. This is
faithful to master, but it is why \texttt{inductiveReduceRecCore.WF} must take saturation as an
external hypothesis instead of deriving it.
\item \textbf{Low.} \srcloc{Lean4Lean/TypeChecker.lean}{946}: \texttt{TypeChecker.etaExpand} has
no call site anywhere under \texttt{Lean4Lean/}. If it is meant to realize the thesis's
$\mathrm{rec}$-normal-form preprocessing (unique.tex \S4, sec:kappa), that intent is undocumented,
and the executable kernel never enforces $\mathrm{rec}$-normal form.
\item \textbf{Low.} \srcloc{Lean4Lean/ForEachExprV.lean}{25}: \texttt{Expr.forEachV} and
\texttt{forEachV'} appear to be dead code; the module is imported by \texttt{TypeChecker.lean} but
neither function has a call site in the kernel files.
\item \textbf{Low.} \srcloc{Lean4Lean/TypeChecker.lean}{593}: \texttt{quickIsDefEq} destructures
\texttt{TypeChecker.State} positionally as \texttt{.mk a1 a2 a3 a4 a5 a6 a7 (eqvManager := m)}.
Adding or reordering a field silently changes what this matches, or fails to compile confusingly;
a \texttt{modifyGet} on the named field would be equivalent and robust.
\item \textbf{Low.} \srcloc{Lean4Lean/TypeChecker.lean}{299}: the \texttt{eagerReduce} marker (a
syntactic \texttt{isAppOfArity} test that switches the definitional-equality strategy, consulted
again at lines 541, 782 and 845) mirrors core Lean's own gadget, so it is correctly absent from
\texttt{divergences.md}. It is nevertheless the one reduction knob in the file with no thesis
counterpart and no local docstring stating what the four sites jointly guarantee.
\item \textbf{Low.} \srcloc{Lean4Lean/Inductive/Add.lean}{470}: \texttt{AddInductive.run} computes
its \texttt{isUnsafe} binding twice (lines 453 and 470), the second shadowing
the first; dead duplication.
\item \textbf{Low.} \srcloc{Lean4Lean/TypeChecker.lean}{493}: \texttt{reducePowMaxExp} silently
declines to reduce \texttt{Nat.pow} at exponents above $2^{24}$. The C++ kernel has the same cap,
but the incompleteness is not listed in \texttt{divergences.md}.
\item \textbf{Low.} \srcloc{Lean4Lean/TypeChecker.lean}{845}: the \texttt{s.isConstOf ``true}
shortcut in \texttt{isDefEqCore'} is an asymmetric special case with neither a thesis counterpart
nor a \texttt{divergences.md} entry.
\item \textbf{Low.} \srcloc{Lean4Lean/TypeChecker.lean}{355} and
\srcloc{Lean4Lean/TypeChecker.lean}{968}: \texttt{reduceProj} is written in an awkward explicit
bind style for the benefit of its proof, and a commented-out \texttt{example}/\texttt{run\_tac}
scratch block is left at the end of the file.
\item \textbf{Low.} \srcloc{Main.lean}{29}: \texttt{FuelConfig.toObj} uses \texttt{panic!} when the
derived JSON encoding is not an object, inside a command line argument parser.
\item \textbf{Informational.} \srcloc{Lean4Lean/PtrEq.lean}{17}: the two pointer-equality
soundness axioms are genuine trusted assumptions of the project. They are honestly documented and
used in only two places, but at \srcloc{Lean4Lean/TypeChecker.lean}{739}
\texttt{ptrEqConstantInfo} gates a different algorithm rather than a shortcut.
\item \textbf{Informational.} \srcloc{Lean4Lean/Verify/Environment.lean}{208}:
\texttt{addDecl.WF}'s \texttt{inductDecl} case is \texttt{sorry}, so everything the \texttt{iota}
branch proves about inductive blocks lives on the model side of that gap. This is master's hole,
not a contributed one, but it bounds what the branch can claim.
\end{itemize}
